\subsection{Wärmeleitungsgleichung}
\begin{emphbox}
    Sei $\Omega\subset\mathbb R^n$ ein Körper mit Massendichte $\rho(x)$, spezifischer Wärme $c(x)$, Wärmeleitfähigkeit $\kappa(x)$, Wärmezufuhr ($f>0$)/-abfuhr ($f<0$) $f(x,t)$
    \begin{equation*}
        c(x)\rho(x)u_t(x,y) - \div(\kappa(x)\nabla u(x,t)) = f(x,t)
    \end{equation*}
\end{emphbox}

\subsection{Randbedingungen}
\begin{sectionbox}
    \subsubsection{Dirichlet-Problem}
    Dirichlet-Randbedingungen legen die Temperatur am Rand des Körpers fest. Gegeben $u^D:\partial\Omega \times [0,T]\rightarrow\mathbb R$.
    Gesucht eine Lösung, so dass \begin{equation*}
        u(x,t) = u^D(x,t) \quad \forall (x,t) \in \partial\Omega\times [0,T]
    \end{equation*}

    \subsubsection{Neumann-Problem}
    Neumann-Randbedingungen schreiben den Wärmefluss nach außen vor. Gegeben $g^N:\partial\Omega\times [0,T]\rightarrow\mathbb R$.
    Gesucht eine Lösung, so dass \begin{equation*}
        \kappa(x)\frac{\partial}{\partial n}u(x,t) = g^N(x,t) \quad \forall (x,t)\in\partial\Omega \times [0,T]
    \end{equation*}
\end{sectionbox}

\subsection{Poisson-Gleichung}